% P55 | 7 Conclusions | Conclusions: the European call and the two languages
% guide: dev_notes/paragraph_guides/P55_conclusions_the_european_call_and_the_two_langua.txt
The discrete and the continuous language agree in the limit: the scaled random walk converges to Brownian motion at rate $N^{-1/2}$ in the Kolmogorov distance, and the CRR price converges to the Black--Scholes price at rate $1/N$ with a sign that depends on the parity of $N$.
The sign pattern is a property of the kink at the strike and can be removed.
For a vanilla payoff a well-chosen lattice beats Monte Carlo by orders of magnitude in accuracy per unit of time, and variance reduction recovers only a part of the gap.
The strength of Monte Carlo is its generality, not its accuracy.
The Greeks can be obtained from the closed form, from the tree and from Monte Carlo, and the three agree.
The error of a discrete delta hedge falls only like $n^{-1/2}$ in the number of rebalances, and a wrong hedging volatility costs far more than a coarser schedule.
The point to remember is that, for a reasonable method, the volatility that goes into the model matters more than the numerical scheme that processes it.
